% Generated draft: requires mathematical and source review.
\documentclass[11pt]{article}
\usepackage[margin=25mm]{geometry}
\usepackage{fontspec}
\setmainfont{Latin Modern Roman}
\usepackage{amsmath,amssymb,mathtools,unicode-math}
\setmathfont{Latin Modern Math}
\usepackage{enumitem,needspace,xurl,hyperref}
\setlength{\parindent}{0pt}
\setlength{\parskip}{5pt}
\setlength{\emergencystretch}{3em}
\newcommand{\N}{\mathbb{N}}
\newcommand{\Z}{\mathbb{Z}}
\newcommand{\Q}{\mathbb{Q}}
\newcommand{\R}{\mathbb{R}}
\newcommand{\C}{\mathbb{C}}
\newcommand{\F}{\mathbb{F}}
\newcommand{\A}{\mathbb{A}}
\newcommand{\PP}{\mathbb{P}}
\newcommand{\E}{\mathbb{E}}
\newcommand{\eps}{\varepsilon}
\newcommand{\dd}{\,\mathrm{d}}
\begin{document}
\section*{UnsolvedMath ID: 30000407}
\textbf{Source problem code:} catalogue entry
\subsection{Definitions and mathematical statement}
Let $p$ be a prime and let $F$ be a field with $\operatorname{char}F\ne p$ containing all $p$-th roots of unity. Fix a separable closure $F^{\mathrm{sep}}$, and let $F(p)$ be the compositum therein of the finite Galois extensions of $F$ of $p$-power degree. Write
\[
 G_F(p)=\operatorname{Gal}(F(p)/F).
\]
This is the maximal pro-$p$ Galois group, rather than an assumption that the full absolute Galois group is pro-$p$. Suppose $G_F(p)$ is topologically finitely generated: some finite subset generates a dense subgroup.

For a valuation $v$ of $F$ and an extension $w$ to $F(p)$, define its decomposition subgroup and inertia subgroup by
\[
 D_w=\{\sigma\in G_F(p):w\circ\sigma=w\},
 \qquad
 I_w=\ker\!\left(D_w\longrightarrow
           \operatorname{Aut}(\kappa(w)/\kappa(v))\right),
\]
using the residue-field action. Different extensions produce conjugate subgroups in the usual valuation-theoretic convention. The source does not state any further restrictions on the valuations.

The conjecture asks whether there exists a free pro-$p$ product decomposition
\[
 G_F(p)\cong\mathop{\ast}\limits_{j\in J}^{\mathrm{pro}\text{-}p}H_j,
\]
where each factor is isomorphic to one of
\[
 \Z_p,\qquad \Z/2\Z\ \text{(allowed only when }p=2\text{)},
 \qquad D_w\ \text{for a valuation with }I_w\ne1.
\]
Here $\Z_p$ is the additive group of $p$-adic integers and the free product is the coproduct in the category of pro-$p$ groups and continuous homomorphisms, not the abstract discrete free product. The statement concerns topological group isomorphism to the allowed factors; it does not require a specified internal embedding of every factor. An empty-product convention would accommodate a trivial group, but the source does not discuss that degenerate case.

The range of valuation types, residue characteristics, and whether the decomposition groups are meant precisely inside $G_F(p)$ for valuations of $F$ must be confirmed in the report. No rank-one, henselian, tame, or residue-characteristic restriction is silently introduced, and no number or uniqueness of factors is claimed.
\subsection{Short English statement}
The conjecture asks whether a finitely generated maximal pro-$p$ Galois group over a field containing the $p$-th roots of unity splits as a free pro-$p$ product of procyclic factors, order-two factors when $p=2$, and valuation decomposition groups with nontrivial inertia.
\subsection{Sources}
\begin{itemize}
\item Local UnsolvedMath registry, permanent ID 30000407, code catalogue entry, ``Structure of Finitely Generated Maximal Pro-$p$ Galois Groups.'' Proposed year 2006; proposer unspecified.
\item \emph{The Arithmetic of Fields} (2006), Oberwolfach Reports, pp.\ 317--378, DOI 10.4171/owr/2006/06; \url{https://doi.org/10.4171/owr/2006/06}.
\item Statement status corrected\_verified; the supplied note says every field, group, and decomposition condition was preserved and only raw TeX delimiters normalized.
\item Literature triage checked 2026-08-21: open. It reports known cases for several field classes but no general resolution, without naming those classes or additional references. Its phrase ``absolute pro-$p$ Galois groups'' is interpreted according to the explicit maximal-pro-$p$ statement. Created 2026-08-25; updated 2026-09-16. Current status and cited claims require independent verification.
\item Registry curation by Ulam AI / UnsolvedMath Contributors, CC BY 4.0; the underlying report retains its supplied provenance. Unreviewed statement-only draft.
\end{itemize}
\end{document}
